\part{Cut-rank laws for graph states, forests, and random width}

\medskip\noindent\emph{Added to the deposit: 30 August 2026.}\par\medskip

For a finite simple graph $G=(V,E)$ and $S\subseteq V$, write
\[
 \rho_G(S)=\rank_{\mathbb F_2}A_G[S,V\setminus S],\qquad
 Z_G(x,y)=\sum_{S\subseteq V}x^{|S|}y^{\rho_G(S)}.
\]
For an edge-labelled graph over $\mathbb F_q$, the same symbol denotes rank
over $\mathbb F_q$, with the standard finite-field Pauli convention for the
associated $q$-level stabilizer graph state.  We write $\operatorname{rw}$
and $\operatorname{bw}$ for rank-width and branch-width, $\nu$ for maximum
matching number, and $\operatorname{Odd}_G(A)$ for the vertices having an odd
number of neighbours in $A$.

The fifteen conjectures below form four connected programmes: quenched and
extreme cut-rank laws for dense random graph states; shape, stability,
extremal, and reconstruction laws for forests; extensive width constants in
random regular, sparse, and percolated graphs; and typical rigidity and orbit
entropy under local complementation.  The graph-state Schmidt-rank identity,
local-complementation invariance, the Oum--Seymour theory of rank-width, and
the known order-of-magnitude bounds for random graphs are calibrations, not
new claims~\cite{HeinEisertBriegel04,Oum05,OumSeymour06,LeeLeeOum12}.

\begin{proposition}[Cut-rank calibrations]\label{as:basic}
\textup{(i)} Across the bipartition $S\mid V\setminus S$, the graph state
$|G\rangle$ has Schmidt rank $q^{\rho_G(S)}$ and entropy
$\rho_G(S)\log q$.  Local complementation preserves the whole labelled
function $S\mapsto\rho_G(S)$.  \textup{(ii)} If $F$ is a forest and a
uniform random bipartition is chosen, its crossing-edge set is a uniform
subset $B\subseteq E(F)$, every $B$ has the same colouring-fibre size, and
\[
 \rho_F(S)=\nu((V(F),B)).
\]
\textup{(iii)} Put $P_{q,0}=1$,
$P_{q,k}=\prod_{i=1}^k(1-q^{-i})$, and
$P_{q,\infty}=\prod_{i\ge1}(1-q^{-i})$.  The limiting nullity law of a
uniform $m\times(m+s)$ matrix over $\mathbb F_q$, for fixed $s\ge0$, is
\[
 \pi_{q,s}(d)=q^{-d(d+s)}
 \frac{P_{q,\infty}}{P_{q,d}P_{q,d+s}},\qquad d\ge0.
\]
\end{proposition}

\begin{proof}
The graph-state amplitude matrix across a bipartition reduces by local
Clifford row and column operations to $\rho_G(S)$ entangled pairs, proving
\textup{(i)}~\cite{HeinEisertBriegel04}.  On a forest, fixing one vertex
colour makes the map from vertex two-colourings to crossing-edge sets
bijective.  A bipartite forest's adjacency bi-matrix has rank over
$\mathbb F_2$ equal to its term rank: a nonzero maximal minor has a unique
perfect matching, since two perfect matchings would create a cycle.  This is
\textup{(ii)}.  Finally, the exact rank count for finite rectangular matrices,
divided by $q^{m(m+s)}$ and then sent to infinity, gives \textup{(iii)}.
\end{proof}

\begin{conjecture}[Quenched balanced-cut nullity universality]\label{as1}
Fix a prime power $q$ and $s\in\{0,1\}$.  Let $G_{2m+s,q}$ have independent
uniform edge labels in $\mathbb F_q$, and set
\[
 \mu_G(d)=\binom{2m+s}{m}^{-1}
 \#\{S:|S|=m,\ m-\rho_G(S)=d\}.
\]
Then $\mu_G$ converges in probability, in total variation, to
$\pi_{q,s}$.
\end{conjecture}

\paragraph{Significance.}
Every fixed cut matrix is uniform, so Proposition~\ref{as:basic}(iii) gives
the annealed law exactly; the content is self-averaging over the exponentially
many, heavily overlapping cuts of one state.  A proof would turn the random
matrix nullity law into a quenched entanglement spectrum and requires a
dependence-sensitive concentration argument, not independent sampling.

\begin{conjecture}[Dependence-sensitive extreme balanced-cut defect]\label{as2}
In the setting of Conjecture~\ref{as1}, let
$D_{\max}(G)=\max_{|S|=m}(m-\rho_G(S))$.  Then
\[
 \frac{D_{\max}(G_{2m+s,q})}
 {\sqrt{(2m+s)\log_q2}}\ \longrightarrow\ 1
\]
in probability.
\end{conjecture}

\paragraph{Significance.}
The one-cut tail $q^{-d^2+O(d)}$ and $2^{2m+s+o(m)}$ balanced cuts force the
displayed first-moment location.  The lower bound is a genuine de-clustering
problem: overlapping cuts share quadratically many entries, and complementary
cuts are identical when $s=0$.  Binary enumeration at orders $10$--$18$ gave
mean normalized values $0.72$--$0.80$, increasing across that range; this
calibrates the scale but is not treated as independence evidence.

\begin{conjecture}[Poisson law for random graph-state distance]\label{as3}
For $G\sim G(n,1/2)$ put
\[
 f(G)=\min_{\varnothing\ne A\subseteq V}
 \bigl(|A|+|\operatorname{Odd}_G(A)\setminus A|\bigr)
\]
and
\[
 \Lambda_n(\ell)=\sum_{a=1}^{\ell}\binom na
 \Pr\{\operatorname{Bin}(n-a,1/2)\le \ell-a\}.
\]
Uniformly along integer sequences for which $\Lambda_n(\ell)$ remains in a
compact subset of $(0,\infty)$, the number of nonempty $A$ whose displayed
weight is at most $\ell$ converges in total variation to
$\operatorname{Poisson}(\Lambda_n(\ell))$.  Consequently
$\Pr(f(G)>\ell)=e^{-\Lambda_n(\ell)}+o(1)$.
\end{conjecture}

\paragraph{Significance.}
The minimum is the distance of the self-dual additive code represented by the
graph state.  Known first-order estimates and the corresponding random
linear-code extreme law do not settle the symmetric generator
$(A,\Gamma A)$, where $A$, $B$, and $A\triangle B$ create structured
dependencies.  A proof would give a sharp finite-blocklength law for random
graph-state error detection.

\begin{conjecture}[Typical rigidity of the labelled cut-rank function]\label{as4}
For $G\sim G(n,1/2)$, with probability tending to one, every graph $H$ on
the same labelled vertex set satisfying
\[
 \rho_H(S)=\rho_G(S)\qquad(S\subseteq V)
\]
is obtainable from $G$ by local complementations, without permuting the
vertices.
\end{conjecture}

\paragraph{Significance.}
The converse is standard, but universal completeness is false: the smallest
known non-locally-equivalent pair with the same cut-rank function has nine
vertices, and relabellings of the Petersen graph give a classical source of
exceptions~\cite{Oum05}.  Exhaustive enumeration through six vertices found
exact equality between cut-rank fibres and labelled local-complementation
orbits, including all $760$ fibres at order six.  The conjecture says that the
alternative isotropic representations behind structured exceptions have
vanishing density.  Since local-unitary equivalence preserves every cut rank,
the claim would also imply generic LU--LC; no redundant LU--LC conjecture is
stated separately.

\begin{conjecture}[Positive-rank-tail log-concavity]\label{as5}
For every connected graph $G$, let
$c_r(G)=\#\{S\subseteq V(G):\rho_G(S)=r\}$ and
$R=\max_S\rho_G(S)$.  Then
\[
 c_r(G)^2\ge c_{r-1}(G)c_{r+1}(G),\qquad 2\le r<R.
\]
\end{conjecture}

\paragraph{Significance.}
The positive-rank restriction is sharp: a connected $13$-vertex graph has
full profile $(2,28,410,1896,3296,2208,352)$, which is not log-concave at
rank one; conditioning on $|S|$ fails already at order four.  The positive
tail passed all $1{,}252$ graphs on at most seven vertices and $8{,}500$
random graphs of orders $8$--$15$.  A proof would expose a previously unknown
negative-dependence or injection principle for graph-state entanglement
multiplicities.

\begin{conjecture}[Full log-concavity for percolated forests]\label{as6}
For a forest $F$, define
\[
 a_j(F)=\#\{B\subseteq E(F):\nu((V(F),B))=j\}.
\]
Then $(a_0(F),\ldots,a_{\nu(F)}(F))$ is log-concave.
\end{conjecture}

\paragraph{Significance.}
By Proposition~\ref{as:basic}(ii), these coefficients are, up to a constant
colouring fibre, the cut-entanglement law of a forest graph state.  Exact
computation passed all $2{,}287$ trees through $13$ vertices, and independent
review extended the check through $18$.  The boundary is real: real-rootedness
fails at seven vertices for $(1,13,30,20)$, while ultra-log-concavity fails for
the $13$-vertex graph6 tree \texttt{LqD?I?@O??g??@} with profile
$(1,75,432,1060,1360,912,256)$.

\begin{conjecture}[Hurwitz stability of the percolated-matching polynomial]\label{as7}
For every forest $F$, every zero of
\[
 H_F(z)=\sum_{B\subseteq E(F)}z^{\nu((V(F),B))}
\]
has strictly negative real part.
\end{conjecture}

\paragraph{Significance.}
$H_F$ is not the classical matching polynomial: its coefficients count edge
subsets according to the matching number they induce.  Its two-state
rooted-tree dynamic programme gives a forest-specific recursion, distinct
from the Hochster-strand Hurwitz conjecture in Part~VII.  Every tree through
$18$ vertices passed exact root tests.  There can be no uniform root-modulus
margin, since $H_{S_n}(z)=1+(2^{n-1}-1)z$ has a negative root tending to zero.

\begin{conjecture}[Path--star stochastic extremality under bond noise]\label{as8}
If $T_p$ retains each edge of an $n$-vertex tree independently with
probability $p\in(0,1)$, then for every integer $k$,
\[
 \Pr\{\nu((S_n)_p)\ge k\}\le
 \Pr\{\nu(T_p)\ge k\}\le
 \Pr\{\nu((P_n)_p)\ge k\}.
\]
For $n\ge4$, equality at every threshold characterizes the star and path,
respectively, up to isomorphism.
\end{conjecture}

\paragraph{Significance.}
This is a distributional strengthening of deterministic matching extremality
and, by Proposition~\ref{as:basic}(ii), an extremal theorem for noisy forest
graph-state entanglement.  Exact coefficientwise comparison through $15$
vertices proves every tested probability $p$ simultaneously, rather than on
a numerical grid.  A proof should couple the rooted matching recursions or
identify a sequence of tree transformations preserving stochastic order.

\begin{conjecture}[Fixed-leaf broom minimum]\label{as9}
For $2\le L\le n-1$, obtain $B_{n,L}$ from $K_{1,L}$ by subdividing one
pendant edge $n-L-1$ times.  Among all $n$-vertex trees $T$ with exactly $L$
leaves,
\[
 \nu((B_{n,L})_p)\ \le_{\rm st}\ \nu(T_p)
\]
for every $p\in(0,1)$, with equality at all thresholds only for the broom up
to isomorphism.
\end{conjecture}

\paragraph{Significance.}
The endpoint classes $L=2$ and $L=n-1$ are unique; the claim begins at
$3\le L\le n-2$.  Exact coefficientwise checks through $15$ vertices found
neither an inequality failure nor another equality case.  The result would
give the first leaf-constrained stochastic extremal law for the matching
number after bond percolation.

\begin{conjecture}[Typical tree reconstruction from the entanglement enumerator]\label{as10}
Let $T_n$ be a uniform labelled tree on $[n]$.  With probability tending to
one, every tree $T'$ satisfying $Z_{T'}(x,y)=Z_{T_n}(x,y)$ is isomorphic to
$T_n$.
\end{conjecture}

\paragraph{Significance.}
Universal reconstruction is false and is not asserted.  The first collision
occurs at order $15$: the nonisomorphic graph6 trees
\texttt{NpCa?C@?a??@?@O???G} and
\texttt{NpCc?D??G?\_@O???g??} have the same $Z$.  Exhaustive enumeration
found this to be the only collision class among all $7{,}741$ unlabelled
$15$-vertex trees, with none through order $14$.  The claim asks whether this
exceptional sparsity persists asymptotically and would make a bipartite
entanglement enumerator a typical complete invariant for trees.

\begin{conjecture}[Rank-width density of random regular graph states]\label{as11}
For each fixed $d\ge3$, if $G_{n,d}$ is a uniform simple $d$-regular graph
and $n$ tends through values with $dn$ even, there is a constant
$\kappa_d\in(0,1/3]$ such that
\[
 \frac{\operatorname{rw}(G_{n,d})}{n}\longrightarrow\kappa_d
\]
in probability.
\end{conjecture}

\paragraph{Significance.}
Expansion gives linear lower bounds and the universal decomposition bound
gives the upper endpoint, but neither yields convergence.  The proposed
mechanism is a limit for the minimum balanced cut-rank profile of the
configuration model, plus recursive balanced decompositions losing only
$o(n)$.  Exact dynamic programming on small cubic graphs gave mean widths
$2.45,2.60,2.90,3.00$ at orders $8,10,12,14$; these values calibrate scale
only.

\begin{conjecture}[Cubic branch-width density and kernel transfer]\label{as12}
There is a constant $\beta_3>0$ such that for a uniform simple cubic graph
$H_N$,
\[
 \frac{\operatorname{bw}(H_N)}N\longrightarrow\beta_3
\]
in probability.  In addition, let $M_N$ be cubic configuration multigraphs
and obtain specified simple graphs $H_N^*$ by resolving every loop and
parallel edge using an explicit set of $o(N)$ edge deletions and insertions.
Replace all but $o(N)$ edges of $H_N^*$ by paths of length at least two,
allow arbitrary further subdivisions, and attach arbitrary trees at
vertices.  For the resulting simple graphs $S_N$,
\[
 \operatorname{rw}(S_N)=\operatorname{bw}(H_N^*)+o(N).
\]
Consequently, if $\varepsilon\downarrow0$, $\varepsilon^3n\to\infty$, and
$G_n\sim G(n,(1+\varepsilon)/n)$, then
\[
 \frac{\operatorname{rw}(G_n)}{\varepsilon^3n}
 \longrightarrow \frac43\beta_3
\]
in probability.
\end{conjecture}

\paragraph{Significance.}
Oum's incidence-graph theorem gives rank-width within one of branch-width
after subdividing every edge once~\cite{Oum05}; the explicit transfer clause
is the missing robustness statement for unreplaced edges, repairs, repeated
subdivision, and pendant trees.  The factor $4/3$ is the cubic-kernel vertex
density.  Existing weakly supercritical results determine only the
$\Theta(\varepsilon^3n)$ scale~\cite{DoErdeKang24}.

\begin{conjecture}[Rank-width surface tension for off-critical percolated grids]\label{as13}
For fixed $p\in[0,1]\setminus\{1/2\}$, let $Q_L(p)$ be bond percolation on
the $L\times L$ square grid.  There is a deterministic $\tau_{\rm rw}(p)$
such that
\[
 \frac{\operatorname{rw}(Q_L(p))}{L}\longrightarrow\tau_{\rm rw}(p)
\]
in probability, with $\tau_{\rm rw}(p)=0$ for $p<1/2$,
$\tau_{\rm rw}(p)>0$ for $p>1/2$, and $\tau_{\rm rw}(1)=1$.
\end{conjecture}

\paragraph{Significance.}
Below criticality all components are logarithmic; above it, percolated-grid
tree-width is linear~\cite{LiMuller17}, and bounded degree lets topological-
minor inequalities transfer a linear lower bound to rank-width.  At $p=1$
the exact grid rank-width $L-1$ fixes the normalization.  The new content is
existence of a surface-tension constant, plausibly through almost-subadditive
rectangle gluing followed by an $o(L)$ comparison with rank-width; the
critical point is deliberately left open.

\begin{conjecture}[Saturation of local-complementation orbit entropy]\label{as14}
Let $\mathcal O_{\rm LC}(G)$ be the set of labelled graphs obtainable from
$G$ by local complementations.  For $G\sim G(n,1/2)$,
\[
 \frac1n\log_2|\mathcal O_{\rm LC}(G)|\longrightarrow\log_2 3
\]
in probability.
\end{conjecture}

\paragraph{Significance.}
The universal bound $|\mathcal O_{\rm LC}(G)|\le3^n$ supplies the endpoint;
the conjecture says that distinct local-Clifford normal forms collide only
subexponentially often~\cite{BahramgiriBeigi07}.  Seeded exact enumeration
gave mean normalized logarithms $1.285,1.320,1.376,1.383$ at orders
$7,8,9,10$, with sample maxima reaching $1.471$, versus
$\log_2 3=1.585$.  Existing orbit maps reach small orders
~\cite{AdcockEtAl20}; a proof would quantify the generic absence of extra
non-Pauli local-Clifford automorphisms and the exponential search space for
graph-state compilation.

\begin{conjecture}[Rank-width density curve for sparse Erd\H os--R\'enyi graph states]\label{as15}
For every fixed $c>0$ there is a deterministic $\kappa_{\rm ER}(c)$ such that
\[
 \frac{\operatorname{rw}(G(n,c/n))}{n}\longrightarrow
 \kappa_{\rm ER}(c)
\]
in probability, where $\kappa_{\rm ER}(c)=0$ for $c\le1$ and
$\kappa_{\rm ER}(c)>0$ for $c>1$.
\end{conjecture}

\paragraph{Significance.}
The qualitative phases are known: subcritical components have bounded
rank-width, critical components are sublinear, and every fixed supercritical
density has linear rank-width~\cite{LeeLeeOum12,DoErdeKang24}.  The new claim
is convergence to a deterministic extensive constant.  Its proposed
mechanism passes to the two-core and kernel, proves a limit for balanced
cut-rank optimization under the local weak degree law, and shows that
attached trees do not alter the leading width.  Conjecture~\ref{as12}
predicts the boundary asymptotic as $c\downarrow1$ but does not imply the
fixed-$c$ limit.
